\clearpage
\section{Extended record: displaced results, ordinary proofs and corrections}
\label{long251:sec:erdos-251-extended-record}

This section holds mathematics that belongs to the complete record for
Problem~\#251 and sits outside the short note.  Every statement here is
labelled by its evidence class in the same vocabulary the note uses.

\subsection{The totient witness and forward propagation}
\label{long251:sec:xr-totient}

The note states the exact denominator law \eqref{long251:eq:den-law} and uses it
directly.  The classical witness that produced it is worth recording, because
it is the form in which the shift length first appears.

\begin{proposition}[a shift of totient length]\label{long251:xr:totient}
Let $T:\N\to\Q$ satisfy the dyadic tail recurrence with integer digits.  If the
reduced denominator $d$ of $T_N$ is odd, then $\sigma_{\ph(d)}(N)$ is an
integer.
\end{proposition}

\begin{proof}
Since $d$ is odd, $2$ and $d$ are coprime, so Euler's congruence gives
$d\mid2^{\ph(d)}-1$.  Writing $2^{\ph(d)}-1=dk$ and $T_N=u/d$ in lowest terms,
$(2^{\ph(d)}-1)T_N=ku$ is an integer, and the integral-shift criterion
transfers this to the shift.
\end{proof}

\begin{proposition}[propagation]\label{long251:xr:propagate}
If $\sigma_h(N)$ is an integer, then $\sigma_h(N+k)$ is an integer for every
$k\ge0$.
\end{proposition}

\begin{proof}
The shift step identity gives
$\sigma_h(N+1)=2\sigma_h(N)-(g_{N+h+1}-g_{N+1})$, an integer combination of an
integer and two digits.  Induct on $k$.
\end{proof}

\noindent These are the totient shift and the propagation theorems cited in the
note.  Two worked instances.  If $\operatorname{den}T_N=3$ then $\ph(3)=2$ and
$3T_N$ is an integer, so $\sigma_2(N)$ is integral while $\sigma_1(N)$ is not;
if $\operatorname{den}T_N=5$ then $\sigma_4(N)$ is integral.  In the orbit with
every digit zero and $T_0=1/12$, the denominator is $12=2^2\cdot3$, the orbit
reaches $T_2=1/3$ with odd denominator at $s=2$, and $h=\ph(3)=2$ is exactly the
shift length seen to be integral from index $2$ onwards.  The totient is a
witness rather than the classification: the exact criterion is
$\operatorname{den}(T_N)\mid2^{h}-1$, and the least admissible $h$ is the
multiplicative order of $2$ modulo the odd part.

\subsection{The finite truncation criterion}\label{long251:sec:xr-truncation}

The note gives the explicit remainder bound for the actual gaps.  The general
truncation statement behind it, valid for any dominating majorant, is the
following.

\begin{proposition}[finite truncation]\label{long251:xr:truncation}
Let $M(n)\ge g_n$ for every $n$, assume the series below converges, and put
\[
 S_{h,N,L}=\sum_{j=1}^{L}\frac{g_{N+h+j}-g_{N+j}}{2^{\,j}},\qquad
 R_{h,N,L}(M)=\sum_{j>L}\frac{M(N+h+j)+M(N+j)}{2^{\,j}} .
\]
If for every $h\ge1$ and every $N_0$ there are $N\ge N_0$ and $L\ge1$ with
$\operatorname{dist}(S_{h,N,L},\Z)>R_{h,N,L}(M)$, then the universal shift
escape of Problem~\ref{long251:prob:escape} holds.
\end{proposition}

\begin{proof}
The part of $\sum_{j\ge1}(g_{N+h+j}-g_{N+j})2^{-j}$ omitted from $S_{h,N,L}$
has absolute value at most $R_{h,N,L}(M)$, so under the displayed inequality
the full sum lies at positive distance from every integer.
\end{proof}

The truncation is an exact modular small-arc problem.  Writing the integral
dyadic block
\[
 D_{h,N,L}=\sum_{j=1}^{L}2^{\,L-j}(g_{N+h+j}-g_{N+j}),
 \qquad S_{h,N,L}=\frac{D_{h,N,L}}{2^{L}},
\]
one has
\[
 \operatorname{dist}(S_{h,N,L},\Z)
 =2^{-L}\min\bigl\{D_{h,N,L}\bmod2^{L},\;
 2^{L}-(D_{h,N,L}\bmod2^{L})\bigr\},
\]
where $D_{h,N,L}\bmod2^{L}$ is the least nonnegative residue, including when
$D_{h,N,L}<0$.  The finite criterion therefore asks the residue of
$D_{h,N,L}$ to avoid the two arcs of radius $2^{L}R_{h,N,L}(M)$ around $0$
modulo $2^{L}$.  A one-block certificate would be a prime-gap theorem
producing such an avoided arc on a logarithmic block.  This rewriting is
paper-level and carries no Lean declaration.

\subsection{Divisor-hitting shift escape}\label{long251:sec:xr-divisorhit}

\begin{problem}[divisor-hitting shift escape]\label{long251:xr:divisor-hit}
For every $r\ge1$, does some positive multiple of $r$ escape cofinally, in the
sense of \eqref{long251:eq:shift-escape} at shift length $mr$ for some $m\ge1$?
\end{problem}

This is an exact reformulation, not a weaker mathematical target.
If $T_0$ is irrational, every positive shift is nonintegral by the block
identity, so divisor-hitting escape follows with $m=1$.  Conversely,
rationality gives a shift $h\ge1$ that is integral at every sufficiently
late index.  For each positive integer $m$, the identity
\[
 T_{N+mh}-T_N=\sum_{j=0}^{m-1}(T_{N+(j+1)h}-T_{N+jh})
\]
makes every multiple $mh$ eventually integral.  Taking $r=h$ contradicts
divisor-hitting escape.  This ordinary deduction uses the checked
rationality classification and a finite telescope; rearranging the
quantifiers supplies no new estimate for the actual prime gaps.

\subsection{Two ordinary proofs about the free-pair producer}
\label{long251:sec:xr-compression}

Both arguments below are ordinary mathematical deductions.  The first uses
the prime number theorem, while the second consumes the free-pair lattice
proved above.  Neither deduction is claimed to be kernel-checked; the exact
tail recurrence and free-pair statements they use have their separate linked
Lean evidence.

\paragraph{State compression.}  Exchanging the order of summation gives
$\sum_{N\le X}T_N\le(p_{X+1}-p_0)+2T_X$.  The prime number theorem
\cite[Ch.~6]{mv2007} also gives
\begin{equation}\label{long251:eq:tail-small-relative-prime}
 \frac{T_X}{p_X}
 =\sum_{j\ge1}\frac{g_{X+j}}{2^jp_X}\longrightarrow0.
\end{equation}
Indeed, for each fixed $j$ the summand without $2^{-j}$ tends to zero because
$p_{X+j}/p_X\to1$.  The usual two-sided bounds $p_n\asymp n\log n$ dominate
$g_{X+j}/p_X$ by $K(1+j)^2$, independently of $X$; multiplication by $2^{-j}$
therefore permits dominated convergence.  Applying the first-moment bound at
$2X$ and using \eqref{long251:eq:tail-small-relative-prime}, there is an absolute
$K_0$ such that
\[
 \sum_{X<N\le2X}T_N\le K_0X\log X
\]
for all large $X$.  Thus, for each fixed $C>1$, Markov's inequality leaves at
least $(1-1/C)X$ indices in $(X,2X]$ with $T_N\le K_0C\log X$.

Under rationality, take $X$ beyond the point where the reduced denominator is
the fixed odd integer $d$.  On this set $T_N$ lies in $d^{-1}\Z_{\ge0}$ and
takes at most $dK_0C\log X+1$ values.  Pigeonhole therefore gives one value at
$\gg_{C,d}X/\log X$ indices.  This compression uses the growth of the actual
primes.  It is false for a general rational integer-digit orbit: the quadratic
countermodel of Proposition~\ref{long251:res:polynomialcountermodel} has the
strictly increasing tails $T_N=2(N+4)^2$.  Even for the primes, equal-tail
pairs have difference zero and hence do not supply the nonintegral difference
required by the free-pair criterion.  The naive form
$\sum_{N\le X}T_N\le p_{X+1}$ is false.

\paragraph{Redundancy of the gap condition.}  Under the free-pair lattice the
difference $T_M-T_N$ is an integer whenever the modulus divides the offset, and
an integer cannot lie in $(\tfrac12,1)$.  The half-window component of the
reduced event of Theorem~\ref{long251:res:signedwindow} is therefore already a
contradiction on its own, and the gap-mismatch condition adds nothing to it.
That changes which component a proof has to produce.

\subsection{A bounded companion to the recurring-values countermodel}
\label{long251:sec:xr-boundedpolignac}

Theorem~\ref{long251:res:polignacfail} carries a logarithmic growth profile so that its
cumulative sequence matches the primes.  The mechanism is visible in a much
smaller example, which is recorded here because it is checkable by hand.

\begin{proposition}[bounded recurring-values countermodel]\label{long251:xr:boundedpolignac}
Put $U_0=4$ and, for $n\ge1$, $U_n=6$ when $n=k!$ for some $k\ge3$ and $U_n=4$
otherwise, and set $a_n=2U_{n-1}-U_n$ for $n\ge1$.  Then $a_n\in\{2,4,8\}$,
the value $2$ occurs at every index $k!$ and the value $4$ at every index
$2\,k!$, so both recur infinitely often at indices divisible by any fixed
$t\ge1$.  The series $\sum_{n\ge1}a_n2^{-n}$ equals $4$ and every tail
$\sum_{j\ge1}a_{N+j}2^{-j}$ equals the integer $U_N$.
\end{proposition}

\begin{proof}
For $k\ge3$ the index $k!$ is even and at least $6$, and $k!-1$ is odd, so the
predecessor of a spike is an ordinary index; hence $a_{k!}=8-6=2$,
$a_{k!+1}=12-4=8$, and $a_n=8-4=4$ elsewhere.  For $k\ge2$ one has
$k!<2\,k!<(k+1)!$, so $2\,k!$ is not a factorial, and $2\,k!-1$ is odd so it is
not a factorial either; therefore $a_{2k!}=4$.  Since $t\mid k!$ for every
large $k$, both values recur in the residue class $0$ modulo $t$.  The identity
$a_n2^{-n}=U_{n-1}2^{-(n-1)}-U_n2^{-n}$ telescopes to
$\sum_{j=1}^{m}a_{N+j}2^{-j}=U_N-U_{N+m}2^{-m}$, and $U$ is bounded, so the
endpoint vanishes.
\end{proof}

The cumulative sequence of this word grows linearly, which is the one property
that separates it from the primes and the reason the note carries the
logarithmic version instead.  Both refute the same generic implication.

\subsection{The measured densities in full}\label{long251:sec:xr-densities}

The note reports the range of the adjacent-mismatch density.  The per-offset
figures below come from the same scan over the $6\,841\,648$ primes under
$1.2\times10^{8}$, with double-precision tails.

\begin{center}
\small
\begin{tabular}{@{}rrrrrr@{}}
\toprule
$h$ & events & density & $\delta=+2$ & $\delta=-2$ & last event at prime\\
\midrule
1 & 56\,427 & 0.008248 & 28\,022 & 28\,405 & 119\,995\,753\\
2 & 31\,979 & 0.004674 & 15\,742 & 16\,237 & 119\,993\,807\\
3 & 30\,233 & 0.004419 & 15\,093 & 15\,140 & 119\,998\,321\\
4 & 29\,262 & 0.004277 & 14\,606 & 14\,656 & 119\,993\,473\\
\bottomrule
\end{tabular}
\end{center}

Every offset $h\le16$ has events, with minimum density $0.00418$ and maximum
$0.008248$; the two signs are near balanced at every offset.  For $h=1$ the
eight band densities, and the same figures after rescaling by the band mean of
$\log p$, run
\[
\begin{array}{l}
 0.011285,\;0.008951,\;0.008303,\;0.007936,\;0.007651,\;0.007507,\;0.007253,\;0.007094;\\[2pt]
 0.17671,\;0.15058,\;0.14420,\;0.14067,\;0.13766,\;0.13667,\;0.13333,\;0.13148 .
\end{array}
\]
The rescaled drift is $0.744$ at $h=1$ and lies between $0.744$ and $0.8121$
across all offsets.  Theorem~\ref{long251:res:sparse} proves that the limiting density
is zero, so the observed decline is the expected behaviour rather than a
failure of the event.

\subsection{Corrections to earlier records}\label{long251:sec:xr-corrections}

\paragraph{The superseded exclusion.}  An earlier internal record narrated a
denominator exclusion at $10^{602}$.  That figure is correct as written and is
superseded twice over: by the kernel-decided floor $2^{589}>10^{177}$ of
Theorem~\ref{long251:res:denominatorfloor}, which is smaller but kernel-checked, and by
the certified continued-fraction exclusion $2^{39997}>10^{12040}$ of
Theorem~\ref{long251:res:cfexclusion}, which supersedes it in strength.  The receipt
records the decimal digit count $12041$; the certified decimal statement uses
the exponent $12040$.

\paragraph{The named missing input.}  An earlier record named
Hardy-Littlewood $k$-tuple correlation of consecutive gaps at a fixed offset as
the missing input.  With the offset freed by Theorem~\ref{long251:res:freepair} the
route needs, for each modulus $t$, cofinally many congruent pairs with
certified nonintegral tail difference.  A further record proposed that two even
values differing by $2$, each occurring infinitely often inside a common index
residue class, would supply the two-condition form.  No proof of that
implication was ever produced, and Theorem~\ref{long251:res:polignacfail} shows that no
proof exists at the level of integer-digit recurrences.  That entry is
withdrawn.

\paragraph{The formal boundary.}  An earlier reading of the short note placed
the real-tail bridge and the real form of the adjacent-pair consumer on the
paper side.  Both are kernel-checked, at
\lword{Erdos251/RealPrimeGapTail.lean}{48}{realPrimeGapTail_recurrence}{line~48}
and
\lword{Erdos251/RealPrimeGapTail.lean}{99}{realTailShift_not_both_integral_of_small_pair_of_digit_ne}{line~99}
of the real prime-gap tail module.  A second reading placed the kernel-decided
denominator floor and the free-pair criterion outside the source revision the
note pins.  Both modules are present at that revision, and both are linked from
the note.
